Online discussions on platforms like Reddit are a productive testbed for fallacy detection: the texts are short, the topics range across politics, science, and lifestyle, and the rhetorical moves include the full inventory of informal fallacies that classical argumentation theory catalogued long before social media~\cite{sahai2021invisiblewall,walton2008schemes}. The Touch\'{e} 2026 lab at CLEF formalises this as a shared task on a Reddit-derived corpus and provides a structured benchmark for systems that combine fallacy detection with argument-scheme classification~\cite{heinrich2026touche}. The lab continues to be a productive platform for argument-aware NLP and remains an important reference for benchmarking systems on these tasks.
% MANDATORY: Touché overview citation per workshop instructions.
% Swap heinrich2026touche for the canonical key (likely kiesel:2026 or
% heinrich:2026) once the 2026 overview paper is published.
Our submission and this notebook contribute to the lab's overall record as documented in the Touch\'{e} 2026 overview paper~\cite{heinrich2026touche}.

Our participation in the 2026 edition~\cite{heinrich2026touche} covers all three subtasks: binary fallacy detection (ST1), eight-way fallacy classification (ST2), and a four-way argument-scheme classification (ST3) organised along a goal axis (practical vs.\ epistemic) crossed with a basis axis (internal vs.\ external)~\cite{macagno2022argumentation}. Each subtask is evaluated under two tracks, a \emph{base} track that uses only original Reddit fields and an \emph{enhanced} track that additionally exposes organiser-supplied paraphrases. We submit one run per subtask and track, totalling six TIRA runs~\cite{frobe2023tira}. All six are fine-tuned RoBERTa-large encoders~\cite{liu2019roberta} with task-specific heads, a small batch of LLM-generated synthetic contrastive pairs, and  for ST2 and ST3 on the enhanced track, a pseudo-labelled top-up of the published unlabelled pool. We arrived at this configuration after running a wider grid that included DeBERTa-v3-base/large~\cite{he2021debertav3}, ModernBERT-large~\cite{warner2024modernbert}, and four Qwen variants from 3B parameters to 32B~\cite{qwen25,qwen3} fine-tuned with LoRA / QLoRA~\cite{hu2022lora,dettmers2023qlora} and chain-of-thought supervision~\cite{wei2022cot}, and after a 5-fold cross-validation comparison on ST3 against the strongest LLM alternative.

A central consideration for our submission was the structural status of the enhanced-track fields. The lab organisers produced these fields with knowledge of the gold labels, which means a classifier trained on them has access to features that an honest pipeline would have to recover from raw text. This shaped both how we read our highest numbers and how the rest of the paper is organised. We report base- and enhanced-track numbers separately throughout, anchor every cross-comparison on base-track or 5-fold CV behaviour, and devote a subsection of the Discussion (\S\ref{sec:disc-leakage}) to a four-part argument that openly concedes what our diagnostic does and does not establish.

The same posture extends to the negative results that shaped the final pipeline: the DeBERTa-v3-large NaN explosions, the LLM detour that did not beat the encoder under cross-validation, and the practical-external (\textsf{pe}) rare-class ceiling that survived targeted augmentation, threshold calibration, and a basis-axis meta-learner. We keep these in the body rather than in footnotes, because the design space they map out is the most useful contribution this notebook can offer a team considering the same shared task next year.

A central theme of our work this year is the honest reading of small-corpus comparisons, separating what the dev numbers and 5-fold CV evidence can support from what they cannot. The rest of the paper is organised as follows. Section~\ref{sec:background} situates the work inside the fallacy-detection and argument-scheme NLP lineage. Section~\ref{sec:system} describes the pipeline, the backbone selection, and the LLM detour. Section~\ref{sec:results} reports per-subtask numbers, the augmentation ablation, and the backbone comparison. Section~\ref{sec:discussion} discusses the encoder--LLM comparison, the enhanced-track leakage question, the \textsf{pe}-ceiling, and the limitations of the present work, and Section~\ref{sec:conclusion} summarises our conclusions.
